\chapter{Methodology} \label{cap:methodology}

This chapter describes the pipeline executed in Chapter~\ref{cap:experimental-evaluation}, summarized in Figure~\ref{fig:executed-pipeline}. The sections follow the stages named in the figure: Data, Feature Construction, Models Used, Strategies, Experimental Evaluation, and Statistical Analysis. Components of the original design that were not executed (additional news corpora, text-sentiment scorers, the extended filter catalogue, combinatorial purged cross-validation) are not described as methodology here; they are listed briefly in Section~\ref{sec:methodology-scope} and developed as future work in Chapter~\ref{cap:conclusion}. Operational details of data acquisition, storage and engine integration are kept in the repository evidence files rather than in the main methodology text.

% Executed pipeline figure for Chapter 3.
\begin{figure}[htbp]
\centering
\begin{tikzpicture}[
  node distance=0.55cm and 0.45cm,
  pipebox/.style={draw, rounded corners, fill=blue!10, minimum width=2.6cm, minimum height=0.75cm, align=center, font=\footnotesize},
  arrow/.style={-{Stealth[length=2mm]}, thick}
]
\node[pipebox] (data) {Data};
\node[pipebox, right=of data] (features) {Feature\\construction};
\node[pipebox, right=of features] (models) {Models\\used};
\node[pipebox, below=of models] (strategies) {Strategies};
\node[pipebox, left=of strategies] (evaluation) {Experimental\\evaluation};
\node[pipebox, left=of evaluation] (stats) {Statistical\\analysis};

\draw[arrow] (data) -- (features);
\draw[arrow] (features) -- (models);
\draw[arrow] (models) -- (strategies);
\draw[arrow] (strategies) -- (evaluation);
\draw[arrow] (evaluation) -- (stats);
\end{tikzpicture}
\caption{Pipeline effectively executed in the experiments reported in Chapter~\ref{cap:experimental-evaluation}.}
\label{fig:executed-pipeline}
\end{figure}


\section{Research Design and Positioning} \label{sec:positioning} \label{sec:research-design}

The study compares two strategy variants on a common market window, execution model and set of price-derived inputs. The \emph{baseline} uses price-derived features only; the \emph{hybrid} adds a news-derived feature family. The question of Section~\ref{sec:objectives} is whether that addition improves risk-adjusted, cost-adjusted return. The comparison is a paired historical simulation, not a randomized intervention, and it does not assert a causal mechanism by which news moves exchange rates. Table~\ref{tab:variables} lists the variables by role.

\begin{table}[htbp]
\centering
\caption{Experimental variables of the present work.}
\label{tab:variables}
\small
\begin{tabular}{p{3.0cm}p{9.8cm}}
\toprule
\textbf{Role} & \textbf{Variables} \\
\midrule
Independent (manipulated) &
Presence or absence of the news feature family (baseline versus hybrid); in the follow-up studies, the entry rule, the exit horizon and the retraining policy. \\
\midrule
Dependent (measured) &
Net return and equity curve after costs, maximum drawdown, number of trades, hit rate and payoff ratio; prediction quality (log-loss, directional accuracy) where a model is evaluated on its own. \\
\midrule
Controlled (fixed) &
Instrument (EUR/USD), execution model and costs, starting deposit, model family, random seed, feature definitions, and the frozen model artifacts and their hashes. \\
\bottomrule
\end{tabular}
\end{table}

The closest published work is that of \citeonline{zhang2025macroalpha} (Section~\ref{sec:hybrid}), which reports cost-adjusted Sharpe ratios of 5.87 on EUR/USD and 4.65 on USD/JPY under five-fold expanding-window cross-validation on daily features. The present work differs in four respects that matter for any comparison of numbers. It trades at minute resolution rather than daily; it uses GDELT \emph{event} intensity rather than article sentiment; it passes the model output through a deterministic rule chain with a per-decision audit trail instead of a single threshold; and it evaluates on chronological splits with a replayed execution engine rather than on cross-validation folds. The numerical results are therefore not directly comparable, and Chapter~\ref{cap:experimental-evaluation} discusses the comparison on that basis.

\section{Data} \label{sec:data-architecture}

\textbf{Prices.} The price series is EUR/USD tick data from the Dukascopy public historical datafeed\footnote{\url{https://www.dukascopy.com/}}. Ticks are decoded, stored as Apache Parquet, and resampled to one-minute bid/ask bars (\texttt{QuoteBar}) with a tick count. A single venue is used throughout, so that all backtest prices come from one quote stream. USD/JPY was part of the original design but was not run in any reported experiment.

\textbf{Event data.} The news signal is built from GDELT 2.0 coded events \cite{leetaru2013gdelt}, materialized from the public BigQuery tables into daily Parquet partitions. GDELT collection for the target window starts on February 18, 2015, so March 2015 is the first fully covered month. A month enters an experiment only after its completion marker exists. Failed acquisition attempts and retry details are retained only as repository evidence.

\textbf{Why event intensity and not text sentiment.} The news input of the executed pipeline is the daily mean Goldstein score of coded events \cite{goldstein1992scale}, not a sentiment score inferred by a language model. The choice reflects both scope and validity. The Global Knowledge Graph carries extracted metadata (themes, entities, tone, quotations) and not the full text of articles\footnote{GDELT Global Knowledge Graph Codebook V2.1, \url{http://data.gdeltproject.org/documentation/GDELT-Global_Knowledge_Graph_Codebook-V2.1.pdf}.}, and materializing it for the whole window was out of reach within the available time and budget. In addition, the Goldstein mapping is a fixed, published scale, so the feature does not depend on a trained model whose training data would have to be checked against the evaluation years. The index is also a deterministic aggregate of a public table, which keeps every run recomputable. The cost is that the text half of the news family stays empty in all reported runs. Whether article text would add information remains open.

\textbf{Chronological allocation.} Each experiment fixes three non-overlapping spans: a family-fitting span for the LightGBM models, a later calibration span for the logistic combiner, and a held-out simulation span replayed by the engine with frozen models. Rows whose forward label horizon crosses a span boundary are purged. The pilot uses February--June 2015 for fitting, July for calibration and September for simulation; the main protocol uses March--December 2015, January 2016 and March--October 2016 (Section~\ref{subsec:one-year-protocol} of Chapter~\ref{cap:experimental-evaluation}). Observation counts are minute rows with overlapping labels and are not independent replications.

\section{Feature Construction} \label{sec:feature-layer}

Features are computed with strict point-in-time discipline: at time $t$, only data with timestamp $\leq t$ enter a feature, indicators are warmed up before any decision, and the training and engine code paths share one feature-definition function. Three price-derived families and one news-derived family are used.

\subsection{Trend features} \label{subsec:ta-structural-features}

The trend family classifies the market direction from exponential moving averages of the minute close (periods 3, 8 and 60) and reads the direction on a higher timeframe through Heikin-Ashi bars \cite{valcu2004heikinashi}. The same perception drives the trend-regime filter F1, which vetoes a decision when the primary and higher-timeframe directions conflict.

\subsection{Indicator features} \label{subsec:indicator-features}

The indicator family contains standard momentum and volatility measures \cite{murphy1999technical}: moving averages, MACD, RSI and ATR. Periods are fixed before any backtest, so that hyperparameter search does not silently extend over the indicator family.

\subsection{Candlestick-pattern features} \label{subsec:pattern-features}

The pattern family encodes candlestick shapes \cite{nison1991candlestick}. Detectors come from the TA-Lib library (\texttt{CDL*} functions)\footnote{\url{https://ta-lib.org/functions/}.} and are emitted only at the bar that closes the pattern. The first pilot used a simple placeholder detector; the TA-Lib detectors enter from the H4 experiments onward (Chapter~\ref{cap:experimental-evaluation}).

\subsection{Event-intensity features} \label{subsec:sentiment-features} \label{subsec:event-features}

The news family is the daily GDELT event intensity of the previous section. Because a daily aggregate is only knowable once its UTC day has ended, it is forward-filled onto the minute grid with a one-day publication lag: the aggregate of day $D$ is visible from 00:00~UTC of day $D+1$. A residual leak is recorded rather than hidden: the aggregate is keyed by the date on which an event occurred, not the date on which GDELT added the record, so an event reported more than one day late is folded into a day whose aggregate is already visible.

\section{Models Used} \label{sec:fusion}

\subsection{Late fusion: per-family models and a logistic combiner} \label{subsec:late-fusion}

One LightGBM classifier is fitted per feature family, using only the family's own features, to predict a fixed-horizon binary label: 1 if the mid price fifteen minute bars after the decision bar is strictly higher than at the decision bar, and 0 otherwise. The baseline has three family models (trend, indicator, pattern) and the hybrid adds a fourth (news).

The meta-learner is a logistic regression with default regularization. \textbf{Its inputs are exclusively the up-move probabilities produced by the family models}, one column per family (three for the baseline, four for the hybrid); raw features such as RSI, MACD or ATR never reach it. Its inputs are therefore on a common probability scale in $[0,1]$. The logistic regression is fitted on the calibration span, on the family models' out-of-sample predictions, so that the in-sample fit of the family models does not leak into the combiner. Its output $\hat{p}_t$ is the single probability read by the terminal filter. The baseline and the hybrid share this specification; the only difference is whether the news-family probability is a combiner input.

Fitted models are saved in a portable, pickle-free format with their provenance (split boundaries, label horizon, input hashes, code revision, library versions, seed 42), and every replay loads the same frozen artifacts and records their hashes.

\section{Strategies} \label{sec:strategies}

\subsection{Deterministic execution as a filter chain} \label{subsec:deterministic-execution}

A strategy turns $\hat{p}_t$ into a trade decision through a chain of deterministic \emph{filters} applied in sequence. Each filter contributes a recommendation and a reason and may veto, short-circuiting the chain to HOLD. Every decision is written to a bar-level decision table with the identifier of the trade it opened, managed or closed, so each trade can be traced to the contribution of every filter. The executed chain has seven filters: F1 trend regime, F2 indicator, F3 candlestick pattern, F4 news context (hybrid only), F5 risk guard (drawdown and exposure caps), F6 capital management (stop, targets, lot from risk) and F7, the terminal rule
\begin{equation*}
\mathrm{decision}_t = \begin{cases}
\text{BUY}  & \text{if } \hat{p}_t > \theta_{\text{high}} \text{ and } v_t = 0 \text{ and } (\neg g \lor r_t = \text{bull}), \\
\text{SELL} & \text{if } \hat{p}_t < \theta_{\text{low}} \text{ and } v_t = 0 \text{ and } (\neg g \lor r_t = \text{bear}), \\
\text{HOLD} & \text{otherwise},
\end{cases}
\end{equation*}
where $r_t$ is the trend regime from F1, $v_t$ the event-veto flag from F4 and $g$ a regime-gate switch. The thresholds $\theta_{\text{high}}$ and $\theta_{\text{low}}$ are fixed in the strategy configuration (0.55 and 0.45 in the main protocol); the calibration span is used to check them and the grid is recorded, but they are not learned. Every filter parameter lives in a YAML descriptor rather than in code, so a variant differs from its base by a configuration diff, and the resolved configuration is saved with each run. The values used are tabulated in Chapter~\ref{cap:experimental-evaluation}.

The separation between a model-based perception layer and this rule-based execution layer is the central design commitment (\emph{agentic perception, deterministic execution}). Model outputs are computed once and read at backtest time, so given the same features and thresholds a replay is reproducible regardless of upstream non-determinism.

\subsection{Strategy variants} \label{subsec:strategy-variants}

The experiments compare the following variants, all on the same chain and execution model. The \textbf{baseline} uses the three price-family models. The \textbf{hybrid} adds the news-family probability to the combiner and the F4 event gate; its comparison with the baseline therefore measures the combined effect of both, and isolating either requires a separate ablation. The follow-up studies add: rule-based news cells (event-intensity trigger on a trailing-quantile threshold, with reversed-sign and always-short controls); an H4/H1/M15 timeframe sweep with TA-Lib patterns and a quote-activity gate; a news-only strategy registered before its outcome; the frozen policy of the adaptive-retraining study; and a screen of a 22-pattern candlestick catalogue, per pattern, in combination and under gated variants. Their definitions and registration dates are summarized with their results in Chapter~\ref{cap:experimental-evaluation}; the full registry is kept in the repository.

\section{Experimental Evaluation} \label{sec:evaluation}

\subsection{Execution engine} \label{sec:execution-engine}

Backtests run on the open-source LEAN engine of QuantConnect, released under the Apache~2.0 license\footnote{\url{https://github.com/QuantConnect/Lean/blob/master/LICENSE}.}, executed locally in a Docker container with its Python algorithm interface. No QuantConnect cloud service is used. The Parquet price store is materialized once into LEAN's native format, and an automated test inside the pinned engine image confirms that training features and the engine's features agree to nine decimal places on every decision bar.

\subsection{Experimental workflow} \label{subsec:workflow-stages}

Each experiment is registered before its evaluation returns are inspected, with hypothesis, variants, label, dates, execution assumptions, primary comparison and candidate-search budget. Family models and combiner are fitted and frozen; the frozen artifacts are replayed over the held-out window under identical assumptions with no fitting in the replay; and the replay is audited (coverage, decision-to-trade joins, zero-trade outcomes) before any metric is read. A registered job is labelled \emph{registered}, \emph{exploratory} or \emph{ad hoc} according to whether its design was written down before any outcome of its window was read. A window that has guided later model or threshold choices becomes development data and is not presented as an untouched test set.

\subsection{Execution model and costs} \label{subsec:execution-assumptions}

The pilot sized orders as a fraction of equity and applied the 0.3-pip median spread of the quote bars. From the confirmation reruns onward the execution model of Section~\ref{subsec:execution-model} is used: stop-first sizing from a fixed risk fraction, one target and a trailing stop, a limit of two concurrent trades, and 1.0 pip of slippage per fill. Both variants inherit identical execution assumptions by construction. Known-answer controls on the engine check costs, sizing, stops, margin and leverage before any outcome is interpreted, because a successful container run does not establish economic realism.

\section{Statistical Analysis of Results} \label{sec:reproducibility}

\subsection{Rule-level significance under multiple comparisons} \label{subsec:rule-significance}

Cross-validation alone does not remove selection bias when many candidates are screened on the same history \cite{aronson2007evidence,bailey2017overfitting}, so the full search history, including unsuccessful variants, accompanies every inferential claim. Inference uses net daily portfolio returns taken from the engine's actual equity, including open-position valuation, at recorded UTC day boundaries; the analyzer neither interpolates missing equity nor substitutes closed-trade returns. The engine's annualized Sharpe ratio is descriptive.

For paired comparisons the estimand is the mean daily difference in net return, tested with the stationary bootstrap of \citeonline{politis1994stationary} on identical circular blocks for both strategies, with a primary and a sensitivity block length fixed before evaluation. In a registered simulation study the procedure passed its size bound only at long windows (twenty-day blocks at twelve hundred observations) and over-rejected at the ninety- and one-hundred-eighty-day windows (null rejection of 9\% to 14\% at the nominal 5\%). Short-window results are therefore diagnostic only, and a mean-return test does not establish Sharpe superiority.

The deflated Sharpe ratio of \citeonline{bailey2014deflated} is implemented from the non-annualized daily Sharpe ratio, skewness and kurtosis, with the selection threshold computed from across-trial Sharpe dispersion and a documented number of effective independent trials; missing selection history yields an unavailable result and not an assumed single trial. For rule cells, a drift control (always-short or always-long on the same bars) separates the effect of the signal from the effect of the currency's path. For pattern screening, a pooled multiple-comparisons check asks whether the best of the screened candidates is distinguishable from the chance outcome of screening that many (Section~\ref{subsec:candles-multiple-comparisons}). Classification-quality metrics are reported separately from, and are never taken as evidence of, trading profitability.

\subsection{Threats to validity} \label{subsec:threats-summary}

The threats are organized along the four biases of \citeonline{hallsmoore2015}. \emph{Optimization bias} is the dominant risk: repeated threshold, gate and strategy selection on the same months means that no reused window is confirmatory, and only windows that no earlier decision had read are treated as out-of-sample. \emph{Look-ahead bias} is addressed by the one-day event lag, the purging of labels at split boundaries, the feature-parity test, and decisions taken at bar close; the residual event-date leak above is disclosed. \emph{Survivorship bias} is not assessed separately, since the study uses one currency pair and not a selected universe of instruments. Registering experiments before viewing outcomes and reporting all variants run limits the room for selective reporting. \emph{Psychological tolerance bias} is not assessed, since the strategies were evaluated by replay and not operated live. The evaluation is a single pair over one year, with one chronological split per experiment, and cannot distinguish a weak regime-dependent edge from the particular year's price path.

\section{Scope: Designed but Not Executed} \label{sec:methodology-scope}

The original design contained more than was run. The following components were not exercised in any reported experiment and are not part of the evidence: additional corpora (FNSPID, CC-News, central-bank releases, Reddit, social-media streams); the dictionary and FinBERT-class sentiment scorers and the GPR index as model inputs; the market-activity feature family; USD/JPY; the extended filter catalogue (F8--F20); the search over filter order and per-filter timeframe; combinatorial purged cross-validation with embargo; rolling refits over a ten-year window; and the four adaptive-retraining policies other than the frozen one. They are developed as future work in Section~\ref{subsec:future-unexecuted} of Chapter~\ref{cap:conclusion}.
